\begin{abstract}
语言智能体在任务自动化方面展现出巨大潜力。随着任务变得日益复杂且执行周期更长，为充分发挥这种潜力，面向多轮任务求解的“子智能体即工具”范式逐渐兴起。然而，现有设计仍缺少对子智能体的\emph{动态抽象}视角，因而适应性不足。为应对这一挑战，我们提出一种统一且与框架无关的智能体抽象，将任意智能体表示为四元组 $\langle Instruction, Context, Tools, Model \rangle$。该四元组可视作能力的组合配方，使系统能够按任务需要创建专用执行器。在此基础上，我们提出智能体系统 \textsc{AOrchestra}。其中央编排器会在每一步实例化该四元组：筛选与任务相关的上下文，选择工具和模型，并通过即时自动创建智能体来委派执行。这一设计能够减少人工工程投入，同时保持框架无关性，并以即插即用的方式支持多种智能体作为任务执行器。它还支持可控的性能与成本权衡，使系统能够逼近帕累托有效前沿。在三个具有挑战性的基准（GAIA、SWE-Bench 和 Terminal-Bench）上，与 Gemini-3-Flash 搭配时，\textsc{AOrchestra} 相比最强基线取得了 $16.28\%$ 的相对提升。代码地址：\url{https://github.com/FoundationAgents/AOrchestra}
\end{abstract}
